\documentclass[11pt, a4paper]{article}
\usepackage[utf8]{inputenc}
\usepackage{amsmath, amssymb}
\usepackage{graphicx}
\usepackage[margin=1in]{geometry}
\usepackage{tcolorbox}
\usepackage{fancyhdr}
\usepackage{hyperref}
\usepackage{listings}
\usepackage{xcolor}

\definecolor{UNIBBlue}{RGB}{0, 51, 102}
\definecolor{codegreen}{rgb}{0,0.6,0}
\definecolor{codegray}{rgb}{0.5,0.5,0.5}
\definecolor{codepurple}{rgb}{0.58,0,0.82}
\definecolor{backcolour}{rgb}{0.95,0.95,0.92}

\lstdefinestyle{mystyle}{
    backgroundcolor=\color{backcolour},   
    commentstyle=\color{codegreen},
    keywordstyle=\color{magenta},
    numberstyle=\tiny\color{codegray},
    stringstyle=\color{codepurple},
    basicstyle=\ttfamily\footnotesize,
    breakatwhitespace=false,         
    breaklines=true,                 
    captionpos=b,                    
    keepspaces=true,                 
    numbers=left,                    
    numbersep=5pt,                  
    showspaces=false,                
    showstringspaces=false,
    showtabs=false,                  
    tabsize=2
}
\lstset{style=mystyle}

\pagestyle{fancy}
\fancyhf{}
\rhead{\textbf{Optimisasi untuk Teknik Elektro}}
\lhead{Lembar Kerja Mahasiswa (LKM) - Minggu 3}
\cfoot{\thepage}

\begin{document}

\begin{center}
    \Large\color{UNIBBlue}\textbf{LEMBAR KERJA MAHASISWA (LKM)}\\
    \large\textbf{Minggu 3: Analisis Solusi LP \& Ekstraksi Output Dual}
\end{center}

\vspace{0.5cm}
\textbf{Nama Praktikan :} \rule{6cm}{0.4pt} \\
\textbf{NPM :} \rule{6.8cm}{0.4pt}

\section*{Tujuan Praktikum}
1. Mahasiswa mampu menggunakan Julia/JuMP atau Python/SciPy untuk memecahkan model LP. \\
2. Mahasiswa mampu mengekstraksi dan menginterpretasikan nilai \textit{dual variable} (shadow price).

\section*{Kasus: Alokasi Frekuensi Komunikasi (Dimensi C3 - C4)}
Sebuah menara BTS melayani dua jenis spektrum layanan: Layanan Voice ($x_1$) dan Layanan Data Broadband ($x_2$). Keuntungan per Mbps alokasi Voice adalah \$40 dan Data Broadband adalah \$50.
Terdapat batasan perangkat keras:
1. Total bandwidth prosesor sinyal (DSP) maksimal 200 Mbps. ($x_1 + 2x_2 \le 200$)
2. Total daya transmisi maksimal 150 unit. ($2x_1 + x_2 \le 150$)
3. Batas regulasi layanan data maksimal 80 Mbps. ($x_2 \le 80$)

Tujuan: Maksimalkan Total Keuntungan $Z = 40x_1 + 50x_2$.

\section*{Aktivitas 1: Implementasi dengan Julia/JuMP (C3)}
\begin{lstlisting}[language=Python]
using JuMP
using HiGHS

model = Model(HiGHS.Optimizer)
@variable(model, x1 >= 0)
@variable(model, x2 >= 0)

@objective(model, Max, 40*x1 + 50*x2)

# Beri nama untuk constraint agar mudah diekstrak nilai dual-nya
@constraint(model, dsp_limit, x1 + 2*x2 <= 200)
@constraint(model, power_limit, 2*x1 + x2 <= 150)
@constraint(model, reg_limit, x2 <= 80)

optimize!(model)

println("Solusi Optimal x1 = ", value(x1))
println("Solusi Optimal x2 = ", value(x2))
println("Keuntungan Maksimal Z = ", objective_value(model))
\end{lstlisting}

\section*{Aktivitas 2: Ekstraksi dan Interpretasi Shadow Price (C4)}
Tambahkan kode berikut pada bagian bawah program di atas:
\begin{lstlisting}[language=Python]
println("Shadow Price Batas DSP = ", dual(dsp_limit))
println("Shadow Price Batas Daya = ", dual(power_limit))
println("Shadow Price Batas Regulasi = ", dual(reg_limit))
\end{lstlisting}

\textbf{Pertanyaan Analisis:}
Berdasarkan nilai \textit{dual variable} (shadow price) yang Anda dapatkan, jelaskan kepada manajemen \textbf{batas sumber daya mana} (DSP, Daya, atau Regulasi) yang harus diprioritaskan untuk di-upgrade kapasitasnya tahun depan agar keuntungan perusahaan meningkat paling besar! Jelaskan alasan kuantitatifnya.

\vspace{2cm}
\textbf{Jawaban:} \\
\rule{\textwidth}{0.4pt} \\
\rule{\textwidth}{0.4pt} \\
\rule{\textwidth}{0.4pt} \\
\rule{\textwidth}{0.4pt}

\end{document}
